\begin{abstract}
Рассматривается расчёт многоканальных систем массового обслуживания с пуассоновскими потоками положительных и отрицательных заявок. Исследуются удаление одной обслуживаемой заявки и мгновенное очищение системы. Для двухфазного гиперэкспоненциального обслуживания описаны матрицы переходов и вычисление стационарных вероятностей методом Такахаси--Таками. Получены выражения для преобразований Лапласа--Стилтьеса времени ожидания, а для модели с очищением --- также времени пребывания. Повторы с новой длительностью описываются переходами между фазами. Для общих прерываний с сохранённой длительностью построена совместная модель, учитывающая первоначальный объём и выполненную часть работы каждой обслуживаемой заявки. Предложена её конечная фазовая аппроксимация с численной проверкой уточнений. В обоих вариантах общих прерываний сохраняется зависимость каналов. Для выборочных прерываний с сохранённой длительностью используется приближение эффективного обслуживания. Проведены проверки по уравнениям баланса, закону Литтла и независимой имитации. Указаны ограничения приближений и условия существования моментов.
\end{abstract}
\textbf{Ключевые слова:} теория массового обслуживания, отрицательные заявки, катастрофы, повторное обслуживание, метод Такахаси--Таками, преобразование Лапласа--Стилтьеса.

\begin{otherlanguage}{english}
\bigskip
\begin{center}\textbf{Calculation of Multiserver Queues with Negative Customers}\end{center}
\textbf{Abstract.} Multiserver queues with independent Poisson streams of positive and negative customers are considered. Negative arrivals either remove one customer in service or clear the entire system. For two-phase hyperexponential service, transition blocks and the Takahashi--Takami stationary calculation are described. Laplace--Stieltjes transforms of waiting times and, for system clearing, sojourn times are obtained. Restart with resampling is represented by phase transitions. For simultaneous fixed-duration restart, a joint model retains each customer's original work requirement and completed work. A finite phase approximation is proposed and checked by numerical refinement. Both simultaneous-restart models preserve dependence between servers. Selective fixed-duration restart uses an effective-service approximation. Balance equations, Little's law and independent simulations are used for verification. The limitations of the approximations and moment-existence conditions are stated.

\textbf{Keywords:} queueing theory, negative customers, catastrophes, service restart, Takahashi--Takami method, Laplace--Stieltjes transform.
\end{otherlanguage}
\end{samepage}

\section{Введение}
Отрицательные заявки описывают воздействия, которые удаляют работу из системы или прерывают её выполнение. Для сетей с положительными и отрицательными заявками известны стационарные распределения в мультипликативной форме при специальных предположениях о потоках, обслуживании и механизме удаления \cite{gelenbe1991,gelenbe1992}. Обзор таких моделей приведён в \cite{artalejo2000}. Для одноканальной системы с произвольным обслуживанием исследованы отрицательные заявки \cite{harrison1996} и мгновенное очищение системы при катастрофе \cite{jain1996}.

В многоканальном случае существенную роль играет распределение остаточных времён обслуживания. Один из подходов к расчёту состоит в фазовой аппроксимации обслуживания и решении уравнений баланса методом Такахаси--Таками \cite{takahashi1976,ryzhikov2018ru}. В настоящей работе этот подход применяется к удалению одной обслуживаемой заявки и к очищению всей системы. Основное внимание уделяется построению переходов и определению временных характеристик. Повторы с новой длительностью учитываются в фазовой модели напрямую. Для общих прерываний с сохранённой длительностью строится совместная модель выполняемых работ; её результаты сопоставляются с приближением независимых эффективных обслуживаний.

\section{Постановка задачи}
Система имеет \(n\) одинаковых каналов и неограниченную очередь. Положительные и отрицательные заявки поступают по независимым пуассоновским потокам интенсивностей \(\lambda>0\) и \(\delta\ge0\). Положительные заявки начинают обслуживание в порядке поступления. Исходные длительности обслуживания \(B\) независимы и одинаково распределены, независимы от входных потоков. Обозначим \(b_k=\mathbb E[B^k]\), \(CV=\sqrt{b_2-b_1^2}/b_1\).

Рассматриваются два механизма воздействия отрицательной заявки:
\begin{enumerate}[leftmargin=1.5em]
\item \textbf{RCS} (Remove Customer in Service): равновероятно выбирается один занятый канал, обслуживаемая заявка удаляется;
\item \textbf{катастрофа}: одновременно удаляются все заявки, находящиеся в обслуживании и в очереди.
\end{enumerate}
В пустой системе отрицательная заявка не оказывает воздействия. При повторном обслуживании удаление заменяется прерыванием: выбранные заявки немедленно начинают обслуживание заново, сохраняя преимущество перед ожидающими. После прерывания либо независимо выбирается новая длительность, либо повторяется первоначальная длительность данной заявки. В обоих случаях выполненная работа теряется; продолжение с места прерывания не рассматривается.

Время ожидания \(W\) отсчитывается от поступления до начала обслуживания, а при очищении --- до начала обслуживания или удаления из очереди. Время пребывания \(V\) заканчивается завершением обслуживания либо удалением. Каждая положительная заявка учитывается один раз. Обозначим \(w_k=\mathbb E[W^k]\), \(v_k=\mathbb E[V^k]\), а через \(q\) --- вероятность завершить обслуживание. Номинальная нагрузка \(\rho=\lambda b_1/n\) используется для задания параметров и не совпадает с фактической занятостью каналов при удалениях. В примерах \(\rho=0{,}7\); для приближённой модели повторного обслуживания дополнительно проверяется \(\lambda\mathbb E[B_{\mathrm{eff}}]<n\).

\section{Базовая фазовая модель и метод Такахаси--Таками}
\subsection{Представление обслуживания и состояния}
Преобразование Лапласа--Стилтьеса (ПЛС) обслуживания аппроксимируется выражением
\[
\beta(s)=\sum_{\ell=1}^{2}y_\ell\frac{\mu_\ell}{\mu_\ell+s},\quad
y_1+y_2=1,\qquad
b_k=k!\sum_{\ell=1}^{2}\frac{y_\ell}{\mu_\ell^k},\quad k=1,2,3.
\]
При \(y_\ell\ge0\) и \(\mu_\ell>0\) это распределение \(H_2\): фаза выбирается при начале обслуживания и не меняется до его завершения. Если согласование моментов приводит к комплексным параметрам или весам вне \([0,1]\), те же формулы используются алгебраически. Такие коэффициенты не интерпретируются как вероятности и интенсивности переходов цепи Маркова; точность аппроксимации требует отдельной проверки. Совпадение первых трёх моментов само по себе не гарантирует точности распределения или его хвоста.

Для вероятностной фазовой модели состояние задаётся парой \((i,j)\): \(i\) --- число заявок в системе, \(j\) --- число обслуживаемых заявок в фазе 2. Число занятых каналов равно \(m=\min(i,n)\), поэтому \(0\le j\le m\). Вероятность состояния обозначим через \(\pi_{i,j}\), вектор вероятностей уровня --- через \(\boldsymbol\pi_i\). Вероятность уровня равна \(p_i=\sum_j\pi_{i,j}\). Расчёт ограничивается уровнями \(0,\ldots,N-1\); \(N\) обозначает число учитываемых уровней.

\subsection{Блоки переходов}
Матрицы \(A_i\), \(B_i\), \(C_i\) задают переходы с уровня \(i\) соответственно на уровни \(i+1\), \(i-1\) и внутри уровня. Диагональная матрица \(D_i\) содержит суммарные интенсивности выхода. В базовой системе \(C_i=0\). При \(i<n\) положительный приход даёт переходы
\[
(i,j)\to(i+1,j):\ \lambda y_1,\qquad
(i,j)\to(i+1,j+1):\ \lambda y_2.
\]
При \(i\ge n\) поступление в очередь не меняет фазы занятых каналов: \((i,j)\to(i+1,j)\) с интенсивностью \(\lambda\).

На уровнях \(1\le i\le n\) завершение обслуживания задаёт переходы
\[
(i,j)\to(i-1,j):\ (i-j)\mu_1,\qquad
(i,j)\to(i-1,j-1):\ j\mu_2.
\]
При \(i>n\) освободившийся канал сразу занимает заявка из очереди, поэтому ненулевые элементы строки \(j\) матрицы \(B_i\) имеют вид
\begin{align*}
B_i(j,j)&=(n-j)\mu_1y_1+j\mu_2y_2,\\
B_i(j,j+1)&=(n-j)\mu_1y_2,\qquad
B_i(j,j-1)=j\mu_2y_1.
\end{align*}
Элементы с индексами вне диапазона \(0,\ldots,n\) отсутствуют. Диагональ определяется суммой выходящих интенсивностей:
\[
D_i(j,j)=\lambda+(m-j)\mu_1+j\mu_2,\quad m=\min(i,n).
\]
Например, при \(n\ge2\):
\[
A_1=\begin{pmatrix}\lambda y_1&\lambda y_2&0\\0&\lambda y_1&\lambda y_2\end{pmatrix},\qquad
B_2=\begin{pmatrix}2\mu_1&0\\\mu_2&\mu_1\\0&2\mu_2\end{pmatrix},
\]
\[
D_2=\begin{pmatrix}\lambda+2\mu_1&0&0\\0&\lambda+\mu_1+\mu_2&0\\0&0&\lambda+2\mu_2\end{pmatrix}.
\]

\subsection{Стационарные вероятности и контроль расчёта}
Внутренние уровни удовлетворяют уравнениям
\[
\boldsymbol\pi_{i-1}A_{i-1}+\boldsymbol\pi_i(C_i-D_i)+\boldsymbol\pi_{i+1}B_{i+1}=0,
\qquad \sum_i\boldsymbol\pi_i\mathbf1=1.
\]
Итерации Такахаси--Таками определяют нормированные векторы уровней \(\boldsymbol t_i\) и отношения \(x_i=p_{i+1}/p_i\). Используются матрицы \((D_i-C_i)^{-1}\) и произведения с соседними блоками; после сходимости вероятности восстанавливаются рекуррентно и нормируются \cite{takahashi1976,ryzhikov2018ru}.

Сходимость итераций и достаточность усечения проверяются раздельно. Контролируются изменение итерационных величин, невязка уравнений баланса, нормировка, а для вероятностной модели --- неотрицательность вероятностей. Затем расчёт повторяется с увеличенным \(N\) и сравниваются целевые характеристики. Малая вероятность последнего уровня сама по себе недостаточна для контроля высоких моментов. При прямых операциях с плотными блоками затраты на проход по уровням имеют порядок \(O(N(n+1)^3)\), объём хранения --- \(O(N(n+1)^2)\).

Для обычной системы \(M/G/n\) с указанной дисциплиной очереди моменты ожидания связаны с факториальными моментами числа ожидающих заявок:
\[
w_k\approx\frac{1}{\lambda^k}\sum_{j=k}^{N-1-n}j(j-1)\cdots(j-k+1)p_{n+j},\quad k\ge1.
\]
При бесконечной сумме равенство точное, если соответствующий момент конечен. Для удаления и очищения временные характеристики вычисляются по ПЛС, описанным ниже. Независимую проверку первых моментов даёт закон Литтла:
\[
\lambda w_1=\sum_i\max(i-n,0)p_i,\qquad
\lambda v_1=\sum_i i p_i,
\]
где учитываются все поступившие заявки, включая удалённые.

\section{Модели отрицательных воздействий}
\subsection{Удаление одной обслуживаемой заявки}
При \(i>0\) каждый из \(m=\min(i,n)\) занятых каналов выбирается с вероятностью \(1/m\). Поэтому блок \(B_i\) получается из базового блока заменой \(\mu_\ell\) на \(\mu_\ell+\delta/m\) в интенсивностях завершения обслуживания. Если очередь непуста, удалённую заявку немедленно заменяет следующая. В частности,
\[
B_2=\begin{pmatrix}2\mu_1+\delta&0\\\mu_2+\delta/2&\mu_1+\delta/2\\0&2\mu_2+\delta\end{pmatrix}.
\]
Для непустых уровней \(D_i^{\mathrm{RCS}}=D_i^{\mathrm{base}}+\delta I\), а \(C_i=0\). В пустом состоянии диагональ остаётся равной \(\lambda\). Переходы показаны на рисунке~\ref{fig:rcs_transitions}. Из баланса потока положительных заявок следует
\[
q=1-\frac{\delta}{\lambda}(1-p_0).
\]

\subsection{Мгновенное очищение системы}
В физической модели катастрофа переводит любое непустое состояние непосредственно в \((0,0)\) с интенсивностью \(\delta\). Чтобы сохранить переходы только между соседними уровнями, вводятся вспомогательные состояния \(\mathcal{D}_i\), \(i\ge1\). При катастрофе из обычного состояния уровня \(i\ge2\) выполняется переход в \(\mathcal{D}_{i-1}\), а из уровня 1 --- непосредственно в \((0,0)\). Далее \(\mathcal{D}_i\to \mathcal{D}_{i-1}\) с интенсивностью \(\gamma>0\), где \(\mathcal{D}_0=(0,0)\). Положительные заявки во вспомогательных состояниях не поступают.

Схема переходов приведена на рисунке~\ref{fig:disaster_transitions}. Первое состояние каждого ненулевого уровня --- \(\mathcal{D}_i\); за ним следуют обычные фазовые состояния. При таком порядке
\[
B_1=\begin{pmatrix}\gamma\\\mu_1+\delta\\\mu_2+\delta\end{pmatrix},\qquad
C_1=\begin{pmatrix}0&0&0\\0&0&0\\0&0&0\end{pmatrix}.
\]
Для \(i\ge2\) первый столбец \(B_i\) содержит \(\gamma\) в строке \(\mathcal{D}_i\) и \(\delta\) во всех обычных строках; оставшаяся часть совпадает с базовым блоком обслуживания. Все \(C_i\) равны нулю. Диагональный элемент вспомогательного состояния равен \(\gamma\), обычного состояния --- базовой интенсивности выхода плюс \(\delta\).

Вспомогательные состояния служат только для решения уравнений. Пусть \(\widetilde\pi\) --- стационарные веса расширенной цепи, а \(c\) --- их сумма по обычным состояниям, включая \((0,0)\). Физические вероятности определяются как
\[
c=\widetilde\pi_{0,0}+\sum_{i\ge1}\sum_j\widetilde\pi_{i,j},\qquad
\pi_{i,j}=\widetilde\pi_{i,j}/c.
\]
Это соответствует удалению времени пребывания во вспомогательных состояниях: при наблюдении только обычных состояний вся цепочка после катастрофы становится мгновенным переходом в пустую систему. Поэтому для вероятностной фазовой модели нормировка восстанавливает исходный процесс при любом \(\gamma>0\). Масса вспомогательных состояний не прибавляется к вероятности пустой системы. Время ожидания и пребывания также вычисляется только по физическим состояниям. Для алгебраической аппроксимации эта операция применяется к алгебраическим весам.

\begin{figure}[!htbp]
\centering
\begin{subfigure}[t]{0.98\linewidth}
\centering
\maybeincludegraphics[width=0.94\linewidth]{figures/figure1a_rcs.png}
\caption{RCS: удаление одной обслуживаемой заявки.}
\label{fig:rcs_transitions}
\end{subfigure}

\begin{subfigure}[t]{0.98\linewidth}
\centering
\maybeincludegraphics[width=0.94\linewidth]{figures/figure1b_disaster.png}
\caption{Очищение: переход в цепочку вспомогательных состояний.}
\label{fig:disaster_transitions}
\end{subfigure}
\caption{Переходы от отрицательных воздействий при \(n=3\). Обозначения внутри обычных состояний показывают числа каналов в фазах 1 и 2. На схеме б стрелка от группы означает переход из каждого её состояния. Переходы от положительных поступлений и завершений обслуживания не показаны.}
\label{fig:negative_transitions}
\end{figure}

\section{Расчёт временных характеристик}\label{sec:times}
\subsection{Ожидание при удалении одной заявки}
Поскольку вход пуассоновский и независим от обслуживания, поступающая заявка видит стационарное распределение состояний. Пока помеченная заявка ожидает, все \(n\) каналов заняты. В состоянии с \(j\) каналами в фазе 2 интенсивность освобождения канала равна
\[
R_j=(n-j)\mu_1+j\mu_2+\delta,\qquad a_j(s)=\frac{R_j}{R_j+s}.
\]
Пусть \(P\) --- матрица вероятностей изменения фаз после освобождения канала и начала обслуживания очередной заявки. Она получается нормировкой строк \(B_{n+1}^{\mathrm{RCS}}\). Если перед поступлением помеченной заявки в системе было \(k\ge n\) заявок, до её обслуживания нужны \(k-n+1\) освобождений. Поэтому
\[
W^*(s)=\sum_{k=0}^{n-1}p_k+
\sum_{k=n}^{N-1}\boldsymbol\pi_k
\left[\operatorname{diag}(\boldsymbol a(s))P\right]^{k-n}\boldsymbol a(s).
\]
Вектор \(\boldsymbol a\) здесь является столбцом. Сначала учитывается время пребывания в текущем состоянии, затем переход; этим определяется порядок матриц. Положительные приходы после помеченной заявки на её ожидание не влияют.

Во время обслуживания интенсивность удаления помеченной заявки равна \(\delta/M(t)\), где \(M(t)\) --- текущее число занятых каналов. Замена этой изменяющейся интенсивности постоянной \(r_m=\delta/m\) даёт приближение
\[
S_m^*(s)=\frac{r_m}{s+r_m}+\frac{s}{s+r_m}\beta(s+r_m).
\]
В расчёте принимается \(m=k+1\), если заявка поступила при \(k<n\), и \(m=n\), если она ожидала. При обозначении положительной части ПЛС ожидания через \(W_+^*(s)\) получаем:
\[
V^*(s)\approx\sum_{k=0}^{n-1}p_k S_{k+1}^*(s)+W_+^*(s)S_n^*(s).
\]
Для \(n>1\) это приближение сохраняет число занятых каналов на момент начала обслуживания. Оно не является точным преобразованием времени пребывания. Среднее \(v_1\) можно независимо проверить по закону Литтла; при \(n=1\) интенсивность удаления во время обслуживания постоянна, и указанное ограничение исчезает.

\subsection{Ожидание и пребывание при очищении}
Пусть \(Y\sim\mathrm{Exp}(\delta)\) --- время до ближайшей катастрофы после поступления помеченной заявки. Введём вспомогательную траекторию, которая начинается из того же стационарного состояния и имеет те же реализации обслуживания, но не содержит последующих катастроф. Обозначим через \(W_0\) время до начала обслуживания помеченной заявки на этой траектории, а через \(Z_0=W_0+B\) --- время до завершения обслуживания.

До ближайшей катастрофы исходная и вспомогательная траектории совпадают. Поэтому непосредственно из механизма очищения следуют равенства
\[
W=\min(W_0,Y),\qquad V=\min(Z_0,Y).
\]
Величина \(Y\) независима от \(W_0\) и \(Z_0\). ПЛС \(W_0^*\) вычисляется по формуле предыдущего подраздела с физическими стационарными весами модели с катастрофами, но без \(\delta\) в интенсивностях освобождения каналов и в матрице \(P\). При этом \(Z_0^*(s)=W_0^*(s)\beta(s)\).

Для независимых \(T\ge0\) и \(Y\sim\mathrm{Exp}(\delta)\) интегрирование функции выживания даёт
\[
\mathbb E[e^{-s\min(T,Y)}]=\frac{\delta}{s+\delta}+\frac{s}{s+\delta}T^*(s+\delta).
\]
Следовательно,
\begin{align*}
W^*(s)&=\frac{\delta}{s+\delta}+\frac{s}{s+\delta}W_0^*(s+\delta),\\
V^*(s)&=\frac{\delta}{s+\delta}+\frac{s}{s+\delta}Z_0^*(s+\delta).
\end{align*}
Вероятность завершения обслуживания равна \(q=Z_0^*(\delta)\). При \(0<q<1\) условные преобразования имеют вид
\begin{align*}
V_{\text{зав}}^*(s)&=Z_0^*(s+\delta)/q,\\
V_{\text{уд}}^*(s)&=\frac{\delta}{s+\delta}\frac{1-Z_0^*(s+\delta)}{1-q}.
\end{align*}
Индексы обозначают завершение обслуживания и удаление. Эти соотношения точны для заданной вероятностной модели; фазовая аппроксимация и усечение пространства состояний вносят отдельные погрешности.

\subsection{Повторное обслуживание с новой длительностью}
После каждого прерывания независимо выбирается новая длительность обслуживания. Для вещественного \(H_2\) достаточно заново выбрать фазу: отсутствие памяти экспоненциального распределения позволяет сохранить состояние \((i,j)\). Число заявок при прерывании не меняется; блоки положительных приходов и завершений обслуживания остаются базовыми. При \(m=\min(i,n)>0\) дополнительные переходы внутри уровня для RCS имеют интенсивности
\[
C_i(j,j+1)=\delta\frac{m-j}{m}y_2,\qquad
C_i(j,j-1)=\delta\frac{j}{m}y_1.
\]
При общем прерывании фазы всех занятых каналов выбираются заново, поэтому
\[
C_i(j,k)=\delta\binom{m}{k}y_2^k y_1^{m-k},\quad k\ne j.
\]
В \(D_i(j,j)\) добавляется сумма этих внедиагональных интенсивностей. Повтор, не изменивший фазовое состояние, не создаёт отдельного перехода. Эта конструкция сохраняет общий момент прерывания и зависимость каналов; усреднение интенсивности RCS не требуется.

Стационарные веса находятся из уравнений баланса. Для временных характеристик дополнительно учитывается помеченная заявка. Пока она ожидает, состояние содержит число заявок перед ней и после неё и фазы занятых каналов; после начала обслуживания --- её фазу, число остальных заявок и фазы остальных занятых каналов. Завершение её обслуживания является поглощающим состоянием. Для переходной части генератора \(\mathcal T\) векторы условных моментов времени до поглощения удовлетворяют
\[
(-\mathcal T)\boldsymbol u_k=k\boldsymbol u_{k-1},\quad
\boldsymbol u_0=\mathbf1,\quad k\ge1.
\]
Усреднение по состоянию, наблюдаемому пуассоновским приходом, даёт моменты пребывания. Аналогичные уравнения с поглощением при начале обслуживания дают моменты ожидания. Для вещественного \(H_2\) это расчёт исходной модели с усечением пространства состояний. При комплексных параметрах или весах вне \([0,1]\) он остаётся алгебраической аппроксимацией.

Независимую проверку даёт формула для одной заявки с пуассоновскими прерываниями постоянной интенсивности \(r\). По уравнению восстановления \cite{pal2017}
\[
B_{\mathrm{eff}}^*(s)=\frac{(s+r)\beta(s+r)}{s+r\beta(s+r)},\qquad
\mathbb E[B_{\mathrm{eff}}]=\frac{1-\beta(r)}{r\beta(r)}.
\]
При \(r=0\) используются исходное распределение и его моменты. При общих прерываниях \(r=\delta\), а будущее обслуживание помеченной заявки независимо от её начального ожидания по свойству отсутствия памяти пуассоновского потока. Поэтому свёртка моментов \(W\) и \(B_{\mathrm{eff}}\) служит дополнительной проверкой моментов \(V\). Независимость обслуживания разных заявок из этого не следует. При RCS интенсивность прерывания во время обслуживания зависит от текущего числа занятых каналов; в расчёте она учитывается через переходы расширенного состояния. В устойчивом режиме без удалений \(q=1\).

\subsection{Повторное обслуживание с сохранённой длительностью}
Если длительность \(B=b\) выбирается один раз и повторяется после каждого прерывания, уравнение восстановления применяется условно по \(b\). Оно даёт
\[
\mathbb E[T\mid B=b]=\frac{e^{rb}-1}{r},\qquad
\mathbb E[T^2\mid B=b]=\frac{2e^{rb}}{r^2}(e^{rb}-1-rb).
\]
Поэтому при \(M_B(t)=\mathbb E[e^{tB}]\)
\[
\mathbb E[T]=\frac{M_B(r)-1}{r},\qquad
\mathbb E[T^2]=\frac{2}{r^2}\bigl(M_B(2r)-M_B(r)-rM'_B(r)\bigr).
\]
Для вещественного \(H_2\) первый момент конечен при \(r<\min_{y_\ell>0}\mu_\ell\), второй --- при \(2r<\min_{y_\ell>0}\mu_\ell\). Для гамма-распределения с масштабом \(\theta\) соответствующие условия имеют вид \(r\theta<1\) и \(2r\theta<1\). Подстановка за полюсом рационального выражения для \(M_B\) не определяет существующий момент.

Для общих прерываний точное состояние процесса содержит число ожидающих \(Q\) и пары \((B_\ell,a_\ell)\) для занятых каналов: \(B_\ell\) --- первоначальная длительность, \(a_\ell\) --- выполненная часть работы. Между событиями \(da_\ell/dt=1\); при \(a_\ell=B_\ell\) обслуживание завершается. Общий переход интенсивности \(\delta\) одновременно заменяет все \(a_\ell\) нулями, сохраняя \(B_\ell\). При наличии очереди завершившую обслуживание заявку заменяет следующая с независимой длительностью. Это марковский процесс с непрерывными координатами и детерминированным движением между скачками. Совместный переход сохраняет зависимость каналов.

Для численного решения закон \(B\) заменяется положительной квадратурой с узлами \(d_h>0\), весами \(\omega_h>0\), \(\sum_{h=1}^J\omega_h=1\). Класс \(h\) выбирается один раз и сохраняется при всех повторах. Условно по этому классу выполнение работы приближается \(K\) последовательными экспоненциальными фазами интенсивности \(K/d_h\). Прерывание возвращает заявку в первую фазу того же класса; общий сброс выполняется одновременно на всех занятых каналах. При конечном \(K\) время прохождения фаз случайно, с дисперсией \(d_h^2/K\); при \(K\to\infty\) оно сходится к \(d_h\). Таким образом, это приближение выполнения сохранённой работы, а не точное описание исходной длительности при конечном числе фаз.

Состояние конечной цепи задаётся длиной очереди и неупорядоченным набором пар «класс, фаза». При всех занятых каналах имеется \(\binom{JK+n-1}{n}\) фазовых состояний. Перестановки одинаковых каналов объединяются без предположения их независимости. Уравнения баланса решаются с разреженными блоками; точность по числу классов, фаз и уровней проверяется раздельно. Моменты ожидания получаются из факториальных моментов \(Q\). Эта связь применима и здесь: пока помеченная заявка ожидает, все каналы заняты, а приходы после неё не влияют на её ожидание.

При общих прерываниях будущее время завершения \(T\) независимо от начального ожидания \(W\). Поэтому для исходного закона \(B\) с \(r=\delta\)
\[
v_1=w_1+\mathbb E[T],\qquad
v_2=w_2+2w_1\mathbb E[T]+\mathbb E[T^2].
\]
Здесь используются приведённые выше точные моменты одной заявки. Для контроля конечной цепи та же свёртка вычисляется с её собственными фазовыми моментами и сопоставляется с законом Литтла для первого момента.

В численных результатах RCS с сохранённой длительностью оставлено отдельное приближение эффективного обслуживания. Сначала находятся первые два момента \(T\) при постоянном \(r\), затем по ним подбирается гамма-распределение, из которого берётся третий момент для двухфазной аппроксимации. Этот дополнительный этап не сохраняет точное распределение \(T\).

Для RCS используется эвристическое замыкание
\[
r=\delta\,\frac{\mathbb E[M]}{\mathbb E[M^2]},\qquad M=\min(i,n),
\]
где распределение \(i\) берётся из вспомогательной системы с эффективным обслуживанием. Это правило не является точным усреднением интенсивности прерывания. Начальное значение \(r=\delta/n\) уточняется итерациями с усреднением старого и нового значений. Допуск равен \(10^{-9}\); при отсутствии сходимости за 100 шагов, бесконечных необходимых моментах или неустойчивости вспомогательной системы результат не выдаётся. Условия существования моментов при постоянном \(r\) относятся к этой вспомогательной модели; они не являются общим критерием для RCS с изменяющейся интенсивностью.

\subsection{Численное получение моментов}
Моменты ПЛС вычисляются по формуле \(\mathbb E[T^k]=(-1)^k(T^*)^{(k)}(0)\). В расчётах удаления и очищения используются центральные конечные разности с девятью узлами и шагом \(h=10^{-3}/b_1\). Для повторов с новой длительностью решаются уравнения первого шага, что исключает численное дифференцирование. Для общих прерываний с сохранённой длительностью используются стационарные вероятности совместной модели и точные моменты одной заявки; для RCS --- приближение эффективного обслуживания. Публикуются первые два момента; их значения сопоставляются с независимыми уравнениями и имитацией. Более высокие моменты без отдельного контроля точности не используются как подтверждённый результат.

\section{Численная проверка}\label{sec:results}
\subsection{Условия эксперимента}
Во всех примерах \(n=3\), \(\lambda=1\), \(b_1=2{,}1\), \(\rho=0{,}7\). Рассматриваются \(CV=0{,}8\) и \(CV=1{,}2\). Для удаления и повторов с новой длительностью \(\delta=0{,}3\), для повторов с сохранённой длительностью \(\delta=0{,}02\). В последнем случае выбор малой интенсивности обеспечивает конечность необходимых моментов вспомогательных распределений в этих примерах.

При \(CV=0{,}8\) имитация использует гамма-распределение со средним \(b_1\) и дисперсией \(b_1^2CV^2\); фазовый расчёт согласует его первые три момента. При \(CV=1{,}2\) и расчёт, и имитация используют вещественное \(H_2\), согласованное с первыми тремя моментами гамма-распределения. При повторе с сохранённой длительностью квадратура и моменты одной заявки строятся для выбранного закона: гамма-распределения в первом случае и \(H_2\) во втором. При повторе с новой длительностью переходы строятся для фазового представления исходного обслуживания.

Для каждой из 12 комбинаций проведено восемь независимых имитационных прогонов. Первые 20\,000 уходов служат прогревом; далее наблюдаются не менее 200\,000 уходов, включая удаления. Генераторы имеют фиксированные различные начальные состояния. Для каждого прогона вычисляются характеристики, затем среднее по прогонам и приближённый 95\%-й доверительный интервал по распределению Стьюдента с семью степенями свободы. Прогрев уменьшает влияние начального пустого состояния, но не заменяет проверку установившегося режима при иных нагрузках.

Расчёт удаления, очищения и RCS с сохранённой длительностью выполнен при \(N=160\) и допуске итераций \(10^{-12}\), затем повторён при \(N=240\). Максимальное абсолютное изменение \(w_1\) и \(v_1\) составило менее \(10^{-12}\). Для вероятностной модели с удалением и очищением дополнительно решены уравнения стационарной цепи с непосредственными физическими переходами, без использования программных блоков Такахаси--Таками. В базовом примере расхождение вероятностей было менее \(3\cdot10^{-10}\), а средних ожиданий с законом Литтла --- менее \(10^{-8}\). Первые два момента ожидания проверены также по независимым уравнениям первого шага. Для очищения проверены среднее пребывание и согласованность условных и безусловных моментов.

При повторе с новой длительностью стационарные уравнения решаются непосредственно разреженным линейным методом. Предел числа заявок после помеченной также увеличивался со 160 до 240; пределы увеличивались как совместно, так и по отдельности. Изменение первых двух моментов ожидания и пребывания не превышало \(2\cdot10^{-8}\). Проверены вырождение в обычную систему при \(\delta=0\), формула Поллачека--Хинчина при \(n=1\) и формула Эрланга при экспоненциальном обслуживании. Для отделения причин расхождений проведены ещё 80 прогонов независимо написанного имитатора: по восемь для десяти сочетаний условий, с прогревом 50\,000 и наблюдением 500\,000 завершений. Среди них сопоставлены общие прерывания и независимые эффективные обслуживания с тем же законом завершения одной заявки. Основные таблицы сохраняют исходные 96 имитационных оценок; дополнительный эксперимент служит отдельной проверкой.

В совместной модели сохранённой работы проверены \(J=3,4\), \(K=2,4,8,12\); при \(CV=1{,}2\) дополнительно использовано \(J=5\), \(K=2,4,8\). Табличные оценки получены при \(J=4\) квадратичной экстраполяцией начальных моментов ожидания по \(1/K\) для \(K=4,8,12\). Коэффициенты определяются сеткой, без использования имитационных значений. Сопоставлены другая тройка фаз \(K=2,4,8\), число классов \(J=3\) и экстраполяция по \(1/\sqrt K\). В последнем сравнении изменение среднего ожидания не превышало 0{,}54\%, стандартного отклонения --- 0{,}94\%. При \(CV=1{,}2\), \(J=4\), \(K=8\) увеличение числа уровней со 100 до 160 изменило \(w_1\) менее чем на \(7\cdot10^{-6}\), \(w_2\) --- менее чем на \(7\cdot10^{-4}\). Разности уточнений характеризуют численную чувствительность; они не являются доказанной границей погрешности стационарных моментов.

\subsection{Результаты и границы точности}
В таблицах~\ref{tab:results}--\ref{tab:sko} приведены расчётные значения и имитационные оценки с полушириной доверительного интервала. Относительное расхождение определяется единообразно: \(\Delta=100(\text{расчёт}-\text{ИМ})/\text{ИМ}\). Это расхождение с оценкой, а не известная точная погрешность относительно стационарного значения.

В восьми примерах удаления и повторного обслуживания с новой длительностью абсолютное расхождение средних ожиданий не превышает 1{,}55\%, средних пребываний --- 0{,}62\%. Для удаления и очищения все приведённые расчётные средние попадают в соответствующие доверительные интервалы имитации. При \(CV=0{,}8\) сравнение включает погрешность фазовой аппроксимации гамма-распределения, а при \(CV=1{,}2\) обслуживание в обеих моделях совпадает. Во всех четырёх примерах повторов с новой длительностью уточнённый расчёт средних и стандартных отклонений также согласуется с доверительными интервалами дополнительной независимой имитации.

\begin{table}[!htbp]
\centering\small\setlength{\tabcolsep}{3pt}
\begin{tabular}{@{}p{4.2cm}cccccc@{}}
\toprule
 & \multicolumn{3}{c}{$w_1$} & \multicolumn{3}{c}{$v_1$} \\
Вариант & Расчёт & ИМ $\pm$ ДИ & $\Delta$, \% & Расчёт & ИМ $\pm$ ДИ & $\Delta$, \% \\
\midrule
\multicolumn{7}{c}{$CV=0,8$} \\
RCS, удаление & 0,4104 & 0,4162 ± 0,0083 & -1,39 & 2,0851 & 2,0907 ± 0,0107 & -0,27 \\
Очищение & 0,1925 & 0,1932 ± 0,0022 & -0,38 & 1,4830 & 1,4835 ± 0,0028 & -0,03 \\
RCS, новая длительность & 1,2067 & 1,1999 ± 0,0181 & 0,57 & 3,4030 & 3,3957 ± 0,0203 & 0,21 \\
Все каналы, новая длительность & 1,8583 & 1,8522 ± 0,0289 & 0,33 & 4,1833 & 4,1759 ± 0,0298 & 0,18 \\
\multicolumn{7}{c}{$CV=1,2$} \\
RCS, удаление & 0,3977 & 0,3998 ± 0,0086 & -0,53 & 1,9223 & 1,9273 ± 0,0107 & -0,26 \\
Очищение & 0,1646 & 0,1642 ± 0,0024 & 0,24 & 1,3117 & 1,3110 ± 0,0036 & 0,05 \\
RCS, новая длительность & 1,0627 & 1,0465 ± 0,0271 & 1,55 & 3,0547 & 3,0361 ± 0,0311 & 0,61 \\
Все каналы, новая длительность & 0,7712 & 0,7722 ± 0,0119 & -0,13 & 2,6625 & 2,6622 ± 0,0109 & 0,01 \\
\bottomrule
\end{tabular}
\caption{Средние ожидание и пребывание при $\delta=0{,}3$. «Новая длительность» означает независимый выбор длительности после каждого прерывания; «все каналы» --- одновременный повтор на всех занятых каналах. ДИ --- полуширина 95\%-го доверительного интервала.}
\label{tab:results}
\end{table}

\begin{table}[!htbp]
\centering\small\setlength{\tabcolsep}{3pt}
\begin{tabular}{@{}p{4.2cm}cccccc@{}}
\toprule
 & \multicolumn{3}{c}{$w_1$} & \multicolumn{3}{c}{$v_1$} \\
Вариант & Расчёт & ИМ $\pm$ ДИ & $\Delta$, \% & Расчёт & ИМ $\pm$ ДИ & $\Delta$, \% \\
\midrule
\multicolumn{7}{c}{$CV=0,8$} \\
RCS, сохранённая длительность & 1,0471 & 1,0323 ± 0,0239 & 1,43 & 3,1751 & 3,1615 ± 0,0265 & 0,43 \\
Все каналы, сохранённая длительность & 1,2711 & 1,2506 ± 0,0242 & 1,64 & 3,4458 & 3,4217 ± 0,0259 & 0,71 \\
\multicolumn{7}{c}{$CV=1,2$} \\
RCS, сохранённая длительность & 1,6107 & 1,6074 ± 0,0365 & 0,21 & 3,7524 & 3,7549 ± 0,0382 & -0,06 \\
Все каналы, сохранённая длительность & 2,2068 & 2,2168 ± 0,0786 & -0,45 & 4,4205 & 4,4312 ± 0,0847 & -0,24 \\
\bottomrule
\end{tabular}
\caption{Средние при повторе с сохранённой длительностью, $\delta=0{,}02$: после каждого прерывания заново выполняется первоначально выбранная работа.}
\label{tab:results_fixed}
\end{table}

\begin{table}[!htbp]
\centering\small\setlength{\tabcolsep}{3pt}
\begin{tabular}{@{}p{4.2cm}cccccc@{}}
\toprule
 & \multicolumn{3}{c}{$\sigma_W$} & \multicolumn{3}{c}{$\sigma_V$} \\
Вариант & Расчёт & ИМ $\pm$ ДИ & $\Delta$, \% & Расчёт & ИМ $\pm$ ДИ & $\Delta$, \% \\
\midrule
\multicolumn{7}{c}{$CV=0,8$} \\
RCS, удаление & 0,9131 & 0,9340 ± 0,0230 & -2,24 & 1,7043 & 1,7093 ± 0,0157 & -0,29 \\
Очищение & 0,5523 & 0,5555 ± 0,0079 & -0,58 & 1,2655 & 1,2652 ± 0,0052 & 0,02 \\
RCS, новая длительность & 1,9620 & 1,9601 ± 0,0489 & 0,09 & 2,6679 & 2,6657 ± 0,0347 & 0,08 \\
Все каналы, новая длительность & 2,7932 & 2,7967 ± 0,0779 & -0,12 & 3,4338 & 3,4350 ± 0,0638 & -0,03 \\
RCS, сохранённая длительность & 1,7708 & 1,7460 ± 0,0447 & 1,42 & 2,4909 & 2,4742 ± 0,0342 & 0,68 \\
Все каналы, сохранённая длительность & 2,1558 & 2,1149 ± 0,0474 & 1,94 & 2,8562 & 2,8203 ± 0,0399 & 1,27 \\
\multicolumn{7}{c}{$CV=1,2$} \\
RCS, удаление & 1,0268 & 1,0252 ± 0,0268 & 0,16 & 2,0976 & 2,0891 ± 0,0159 & 0,40 \\
Очищение & 0,5699 & 0,5696 ± 0,0086 & 0,05 & 1,4337 & 1,4333 ± 0,0040 & 0,03 \\
RCS, новая длительность & 2,0208 & 1,9836 ± 0,0444 & 1,87 & 3,0849 & 3,0572 ± 0,0319 & 0,91 \\
Все каналы, новая длительность & 1,5746 & 1,5817 ± 0,0274 & -0,45 & 2,6559 & 2,6569 ± 0,0181 & -0,04 \\
RCS, сохранённая длительность & 2,7956 & 2,8138 ± 0,0911 & -0,65 & 3,8640 & 3,8838 ± 0,0700 & -0,51 \\
Все каналы, сохранённая длительность & 3,8803 & 3,8867 ± 0,2009 & -0,16 & 4,8643 & 4,8699 ± 0,1703 & -0,11 \\
\bottomrule
\end{tabular}
\caption{Стандартные отклонения ожидания и пребывания. Варианты повторного обслуживания и интенсивности соответствуют таблицам 1 и 2.}
\label{tab:sko}
\end{table}


Раздельная проверка выявляет существенный вклад общих прерываний. При сохранённой длительности и \(CV=1{,}2\) независимый имитатор даёт среднее ожидание \(2{,}229\pm0{,}050\) для общих прерываний и \(2{,}049\pm0{,}034\) для независимых эффективных обслуживаний. Здесь указаны приближённые 95\%-е интервалы по восьми дополнительным прогонам. Распределение завершения одной заявки совпадает, но совместная работа каналов различается. Уточнение отдельных эффективных моментов не устраняет эту зависимость.

Совместная модель сохранённой работы даёт в этом примере \(w_1=2{,}207\) и \(\sigma_W=3{,}880\). Относительно одной и той же независимой имитационной оценки абсолютное расхождение среднего ожидания уменьшилось с 5{,}91\% для приближения эффективного обслуживания до 1{,}01\%, стандартного отклонения --- с 10{,}56\% до 0{,}54\%. При обоих значениях \(CV\) расчётные средние и стандартные отклонения для общих прерываний попадают в интервалы дополнительной имитации. Это подтверждает улучшение для рассмотренных условий, но не устанавливает равномерной точности при иных нагрузках и законах обслуживания. Для RCS с сохранённой длительностью остаются ограничения эвристического усреднения интенсивности и согласования эффективных моментов.

\FloatBarrier
\section{Заключение}
Описан расчёт стационарных вероятностей многоканальной фазовой системы с удалением одной обслуживаемой заявки и с мгновенным очищением. Для очищения предложено использовать вспомогательные состояния с последующим исключением времени пребывания в них. Это сохраняет структуру соседних уровней и даёт физическую нормировку вероятностей. Формулы ожидания учитывают время до последовательных освобождений каналов; формулы пребывания при очищении следуют из совпадения траекторий до ближайшей катастрофы.

Для повторов с новой длительностью построен расчёт по фазовым переходам и уравнениям для помеченной заявки. Он учитывает общие прерывания и изменение индивидуальной интенсивности RCS. Для общих прерываний с сохранённой длительностью построены совместный процесс выполняемых работ и его конечная фазовая аппроксимация. Учёт зависимости каналов уменьшил расхождение с независимой имитацией; остаточная численная чувствительность проверена по классам, фазам и уровням. Для RCS с сохранённой длительностью сохранено приближение эффективного обслуживания. Формулы для одной заявки с постоянной интенсивностью прерываний отделены от приближений многоканальной системы. Программная реализация основана на библиотеке Most-Queue \cite{mostqueue2026}; параметры, начальные состояния генераторов и результаты прогонов сохранены вместе с материалами расчёта.

\printbibliography
\end{document}
